\section{10 September and 18 October 1988: The Liturgical Decree and the Erection}

\subsection{The decree of 10 September: a faculty, and its geography}

Six weeks before the erection, the Commission issued a short Latin decree
granting the liturgical faculty. Its operative sentence is one period,
and every clause of it recurs, verbatim in substance, in the papal decree
of 2022:

\begin{quote}\small
In virtue of the faculty granted to it by the Supreme Pontiff John Paul
II, the Pontifical Commission Ecclesia Dei concedes to that which is
called the ``Fraternity of St.\ Peter'', founded July 18, 1988 and
declared of ``Pontifical Right'' by the Holy See, the faculty of
celebrating Mass, and carrying out the rites of the sacraments and other
sacred acts, as well as fulfilling the Divine Office according to the
typical edition of the liturgical books in force in the year 1962; namely
the Missal, Ritual, Pontifical, and Roman Breviary. This faculty may be
used in their own churches or oratories; otherwise it may only be used
with the consent of the Ordinary of the place, except for the celebration
of private Masses.\endnote{Pontifical Commission \work{Ecclesia Dei},
\latin{Decretum}, 10 September 1988, English text published by the
Fraternity (fssp.org, ``Documents / Foundation''), marked ``[Original:
Latin],'' read 2026-07-25, signed by Augustin cardinal Mayer, president,
and Camille Perl, secretary. The Latin is quoted by Pietras,
art.\ cit., 173--174 n.~12, who cites it \latin{pro manuscripto}: ``facultatem
concedit Missae sacrificium celebrandi, ritus sacramentorum aliosque
sacros ritus peragendi necnon Officium Divinum persolvendi secundum
editiones typicas librorum liturgicorum anno 1962 vigentium, scilicet
Missale, Rituale, Pontificale, et Breviarium Romanum. Qua facultate uti
poterunt in ecclesiis vel oratoriis propriis; alibi vero nonnisi de
consensu Ordinarii loci, excepta Missae privatae celebratione.'' Bounded
negative: this decree was not located in \work{Acta Apostolicae Sedis},
on vatican.va, or in any other official publication.}
\end{quote}

Four observations, each of which matters later.

\emph{The grant is of four books, not of a Mass.} Missal, Ritual,
Pontifical and Breviary: the sacramental rites, the pontifical rites
including ordination, and the Office, as well as Mass. In 2021 the
Congregation for Divine Worship would answer, of the general law, that
the diocesan bishop may permit the 1952 Ritual but not the Pontifical.
The Fraternity's grant of 1988 is wider than what the general law of 2021
allows a bishop to permit, and the difference is the substance of the
Fraternity's later difficulties.

\emph{The grant is territorial in a precise way.} In its own churches and
oratories the faculty operates by itself; elsewhere it requires the
consent of the local Ordinary, save for private Masses. Nothing in the
decree exempts the Fraternity from the diocesan bishop outside its own
places, and nothing in it makes the faculty depend on him inside them.

\emph{The decree calls the body ``of pontifical right'' five weeks before
the erection decree.} This is a genuine oddity of the file. The decree of
10 September describes the Fraternity as ``founded July 18, 1988 and
declared of `Pontifical Right' by the Holy See''; the express erection
in pontifical right is dated 18 October. Either an earlier declaration
existed and has not been published, or the September decree anticipates
the October one, or the phrase is doing loose duty for the readiness
declared on 22 July. No document reached for this account
resolves it, and the resolution is not supplied here.\endnote{Bounded
negative. No declaration of pontifical right earlier than 18 October 1988
was located in \work{Acta Apostolicae Sedis}, on vatican.va, or among the
Fraternity's published founding documents. The Latin of the September
decree quoted by Pietras carries the same anticipation as the English
text, so the difficulty is not an artefact of translation.}

\emph{The decree is a faculty, and faculties are held by their grantor's
title.} That is not a criticism; it is the ordinary structure of such
grants, and it is exactly the structure that the papal decree of 11
February 2022 reproduces.

\subsection{The decree of erection, 18 October 1988}

The erection decree carries the Commission's protocol number 234/88 and
is signed by Mayer and Perl. Its Latin is preserved in the canonical
literature and its English text is published by the Fraternity. It does
five distinct things.\endnote{Pontifical Commission \work{Ecclesia Dei},
\latin{Decretum erectionis ``Fraternitatis Sacerdotalis Sancti Petri''},
18 October 1988, Prot.\ N.~234/88, English text published by the
Fraternity (fssp.org, ``Documents / Foundation''), marked ``[Original:
Latin],'' read 2026-07-25. The Latin is printed in \work{Archiv für
katholisches Kirchenrecht} 157 (1988) 467--468, as cited by Pietras,
art.\ cit., 174--176 nn.~14, 16--18, who quotes the operative clauses; that
periodical itself was not reached for this account, and the Latin used
here is the Latin as quoted by Pietras. Bounded negative: the decree is
not in \work{Acta Apostolicae Sedis} and was not located on vatican.va.}

\emph{One: it erects.} ``This Pontifical Commission \latin{Ecclesia Dei},
in virtue of the special faculties granted to it by the Sovereign
Pontiff, and graciously accepting the petition of the Reverend Father
Josef Bisig, by this selfsame Decree erects the Priestly Fraternity of
St.\ Peter as a clerical society of Apostolic Life with Pontifical Right,
according to the prescribed norms of Canon Law and with all the legal
consequences involved.'' The Latin is \latin{in societatem clericalem
vitae apostolicae iuris pontificii \ldots hoc Decreto iuxta omnes iuris
effectus erigit}. Note the express reference back to the special
faculties: the decree is an exercise of \latin{Quia peculiare munus}.

\emph{Two: it states the purpose in the terms of the papal faculty.} The
Fraternity ``proposes the sanctification of priests through the exercise
of the pastoral ministry, particularly in conforming its life to the Most
Holy Sacrifice of the Mass and by observing the liturgical and
disciplinary traditions invoked by the Roman Pontiff in the Apostolic
Letter Ecclesia Dei of 2 July 1988.'' \latin{Observantiam traditionum
liturgicarum et disciplinarum a Romano Pontifice memoratarum}: the
observance of those traditions is written into the object of the erection
itself.

\emph{Three: it fixes the canonical furniture.} The erection ``brings
with it the rights enumerated in Canon 611''---that is, the rights that
follow from a diocesan bishop's consent to the establishment of a house:
to lead a life according to the institute's proper character and purpose,
to undertake works proper to it, and to have a church and exercise sacred
ministries there, observing the law. The Fraternity ``is regulated by the
norms of the Code of Canon Law, the prescriptions of this Decree, its own
Constitutions and other appropriate laws.''

\emph{Four: it grants the books again, and to a wider class.} ``The
members of the Priestly Fraternity of St.\ Peter, as well as other priests
who are guests in houses of the Fraternity or who exercise the sacred
ministry in their churches, are conceded the use of the liturgical books
in force in 1962.'' The Latin is \latin{usus conceditur librorum
liturgicorum iam anno 1962 vim habentium}. This is broader than the
September decree: it reaches guests and visiting ministers, not only
members.

\emph{Five: it binds the Fraternity into the diocesan structure and
reserves its supervision.} Members are ``with particular diligence to
seek communion with the bishop and diocesan priests according to Canons
679--83''; in pastoral ministry ``the prescriptions of the law are to be
observed, particularly in what concerns the valid and licit celebration
of the Sacraments of Penance and Marriage, as well as what is laid down
in canon 535 concerning the transcription of these events in the parish
registers.'' And, in the clause that reproduces faculty 6 of the papal
rescript: ``the Priestly Fraternity of St.\ Peter is under the authority
of the Sovereign Pontiff as transmitted to this Pontifical Commission for
all that concerns it, until otherwise provided for''---\latin{quoad omnes
effectus subicitur auctoritati Sedis Apostolicae ope huius Pontificiae
Commissionis, donec aliter provideatur}.

Finally, the decree approves the constitutions for three years and names
Father Bisig superior general for three years, and records that John Paul
II, in an audience granted to the undersigned Cardinal President on 18
October 1988, ratified the decree and ordered its publication.

\subsection{What kind of body this is}

It is worth stating plainly, because the term is regularly misused. A
society of apostolic life is not an institute of consecrated life. Canon
731 §\,1 defines such societies as those whose members, without religious
vows, pursue the apostolic purpose proper to the society and, leading a
life in common according to their proper manner of life, strive for the
perfection of charity through the observance of the
constitutions.\endnote{\work{Codex Iuris Canonici} (1983), cann.~731--746,
read in the Latin text of Book II in the Holy See's web presentation
(\url{https://www.vatican.va/archive/cod-iuris-canonici/latin/documents/cic_liberII_la.html})
on 2026-07-25; the exact bytes read match the artifact registered in this
repository's source library. The English rendering used in this section
follows the Holy See's own English page for cann.~731--746
(\url{https://www.vatican.va/archive/cod-iuris-canonici/eng/documents/cic_lib2-cann731-746_en.html}),
read the same day. No project translation is offered.} Canon 732 applies
to societies of apostolic life what canons 578--597 and 606 establish for
institutes of consecrated life, ``with due regard for the nature of each
society''---which is how canon 586 on the just autonomy of life, canon
587 on the fundamental code, and canon 593 on the immediate subjection of
bodies of pontifical right to the Apostolic See in internal governance
and discipline reach the Fraternity at all.

Three consequences follow that the reader should carry through the rest
of this account.

\begin{itemize}
\item \textbf{The members take no religious vows.} They are bound by the
common obligations of clerics and by the bond their own constitutions
establish. Nothing about the Fraternity is monastic.
\item \textbf{``Pontifical right'' is a statement about the supervising
authority, not about exemption.} Canon 593 subjects such a body
immediately and exclusively to the Apostolic See in internal governance
and discipline. Canon 738 §\,2 subjects the members to the diocesan
bishop ``in those matters which regard public worship, the care of souls,
and other works of the apostolate, with attention to cann.\
679--683''---which is exactly the block of canons the erection decree
names.
\item \textbf{A house requires the diocesan bishop's previous written
consent} (canon 733 §\,1), and once erected it carries the canon 611
rights the decree expressly invokes. The Fraternity's growth is
therefore, at every step, a series of episcopal consents, not a series of
Roman grants.
\item \textbf{Clerics are incardinated in the society itself} unless the
constitutions provide otherwise (canon 736 §\,1). The Fraternity's
published statistics distinguish incardinated members from those
incorporated \latin{ad annum}, postulants and associates, and this
account keeps those categories distinct wherever it reports numbers.
\end{itemize}

\actrecord{Decree of erection, Pontifical Commission \work{Ecclesia Dei},
18 October 1988, Prot.\ N.~234/88, ratified in the papal audience of the
same day.}{That the Fraternity was erected as a clerical society of
apostolic life of pontifical right in exercise of the special faculties
of \latin{Quia peculiare munus}; that observance of the liturgical and
disciplinary traditions named in \work{Ecclesia Dei} entered the object
of the erection; that the use of the 1962 books was conceded to members
and to guests and ministers in their churches; that canons 611, 535 and
679--683 were expressly applied; that the constitutions were approved for
three years and Father Bisig named superior general for three years; and
that the Fraternity was subjected to the authority of the Apostolic See
through the Commission \latin{donec aliter provideatur}.}{A decree of a
commission, never printed in the official gazette, known in Latin only
through a canonical periodical not reached for this account and in
English only through the beneficiary's own publication. It settles the
canonical figure and the liturgical grant; it does not state what the
constitutions contained, and it does not say that the 1962 books are the
members' only permitted books.}
